% ****** Start of file apssamp.tex ******
%
%   This file is part of the APS files in the REVTeX 4.2 distribution.
%   Version 4.2a of REVTeX, December 2014
%
%   Copyright (c) 2014 The American Physical Society.
%
%   See the REVTeX 4 README file for restrictions and more information.
%
% TeX'ing this file requires that you have AMS-LaTeX 2.0 installed
% as well as the rest of the prerequisites for REVTeX 4.2
%
% See the REVTeX 4 README file
% It also requires running BibTeX. The commands are as follows:
%
%  1)  latex apssamp.tex
%  2)  bibtex apssamp
%  3)  latex apssamp.tex
%  4)  latex apssamp.tex
%
\documentclass[%
 reprint,
%superscriptaddress,
%groupedaddress,
%unsortedaddress,
%runinaddress,
%frontmatterverbose, 
%preprint,
%preprintnumbers,
%nofootinbib,
%nobibnotes,
%bibnotes,
 amsmath,amssymb,
 aps,
%pra,
%prb,
%rmp,
%prstab,
%prstper,
%floatfix,
]{revtex4-2}

\usepackage{graphicx}% Include figure files
\usepackage{dcolumn}% Align table columns on decimal point
\usepackage{bm}% bold math
\usepackage{multirow}
\usepackage{color}
\usepackage{changepage}
\usepackage{hyperref}
%\usepackage{hyperref}% add hypertext capabilities
%\usepackage[mathlines]{lineno}% Enable numbering of text and display math
%\linenumbers\relax % Commence numbering lines

%\usepackage[showframe,%Uncomment any one of the following lines to test 
%%scale=0.7, marginratio={1:1, 2:3}, ignoreall,% default settings
%%text={7in,10in},centering,
%%margin=1.5in,
%%total={6.5in,8.75in}, top=1.2in, left=0.9in, includefoot,
%%height=10in,a5paper,hmargin={3cm,0.8in},
%]{geometry}

\begin{document}

\preprint{APS/123-QED}

\title{Tidal heating as a direct probe of Strangeness inside Neutron stars}% Force line breaks with \\
%\thanks{A footnote to the article title}%

\author{Suprovo Ghosh}
\email{suprovo@iucaa.in}
\author{Bikram Keshari Pradhan}
%\email{bikramp@iucaa.in}
 %Lines break automatically or can be forced with \\
\author{Debarati Chatterjee}%
% \email{debarati@iucaa.in}
\affiliation{%
 Inter-University Centre for Astronomy and Astrophysics,
Pune University Campus,
Pune 411007, India\\
 %This line break forced with \textbackslash\textbackslash
}%

%\date{\today}% It is always \today, today,
             %  but any date may be explicitly specified

\newcommand{\sg}[1]{\textsf{\color{blue}{ #1}}}
\newcommand{\dc}[1]{\textsf{\color{red}{ #1}}}
\newcommand{\bkp}[1]{\textsf{\color{green}{ #1}}}

\begin{abstract}
The cores of neutron stars (NS) reach densities several times the nuclear saturation density and could contain strangeness containing exotic particles such as hyperons. During the binary inspiral, viscous processes inside the NS matter can damp out the tidal energy induced by the companion and convert this to thermal energy to heat up the star.  We demonstrate that the bulk viscosity originating from the non-leptonic weak interactions involving hyperons is several orders of magnitude higher than the standard neutron matter shear viscosity in the relevant temperature range of $10^6-10^9$K and for heavier mass NSs ($M \geq 1.6M_{\odot}$) that contain a significant fraction of hyperons in their core, the bulk viscosity can heat up the stars upto $0.1 - 1$ MeV before the final merger. This ``tidal heating" process also introduces a net phase shift of $10^{-3}-0.5$ rad, depending on the component mass, in the gravitational wave (GW) signal that can potentially be detected using current and future generation GW detectors. Such a detection would be the direct confirmation of the presence of hyperons inside the NS core, having a great significance for the study of dense matter under extreme condition.
%thus motivating for a more detailed assessment of the tidal heating from bulk viscous dissipation in the gravitational-wave inspiral signal from merging neutron stars. 

\end{abstract}

%\keywords{Suggested keywords}%Use showkeys class option if keyword
                              %display desired
\maketitle

%\tableofcontents

\section{\label{sec:intro} Introduction}
Neutron stars are unique astrophysical dense objects to study  cold  and  dense  matter  under  extreme  conditions far beyond the reach of terrestrial experiments. At such high densities inside the NS core, strangeness containing exotic  particles, such  as hyperons, kaons or even deconfined quark matter, which  are  observed  only  briefly  in  particle  accelerators, can become stable components  due to weak equilibrium and have important consequences for the equation of state (EoS), composition, structure and evolution of these compact objects~\cite{GlendenningBook,Lattimer2004,Blaschke2018}. %\dc{cite some more recent reviews of NSs} 
Earlier various astrophysical observations e.g., the measurement of high mass NS~\cite{Weissenborn2011,Weissenborn_2012}, radius measurement from low-mass X-ray binaries~\cite{Constanca2013,Fortin2015}, thermal evolution~\cite{Anzuini_2021} and GW emission from  hot and rapidly rotating newly born neutron stars due to unstable r-modes~\cite{Lindblom2002,Debi_2006,Jha2022} have been used to constrain the strangeness content inside their interior~\cite{Chatterjee_2016}. The recent observation of the binary neutron star (BNS) merger event GW170817~\cite{Abbott2017,Abbott2018,Abbott2019} using the LIGO–Virgo–Kagra (LVK) global network of GW detectors~\citep{AdvLIGO2015,Abbott2016,AdvVIRGO2014,Kagra} along with detection of its electromagnetic counterpart has opened up the modern era of multi-messenger astronomy~\citep{MultiGW170817}. The measurement of dimensionless tidal deformability~\cite{Hinderer} from the GW170817 merger along with recent NICER estimates of mass-radius posteriors from two pulsars, PSR $J0030+0451$~\cite{NICER0030_Miller,NICER0030_Riley} and PSR $J0740+6620$~\cite{NICER0740_Miller,NICER0740_Riley}, have significantly constrained the unknown neutron star EOS~\cite{Dietrich2020,Pang_2021,Issac2021,Ghosh2022}.  \\

Tidal interaction due to a companion in the binary system  transfer mechanical energy and angular momentum to the star at the expense of the orbital energy and the damping mechanisms inside the star convert this energy to heat which we refer to as ``tidal heating". Tidal heating could potentially cause mass ejection in a radiation-driven outflow prior to the merger~\cite{Rees1992} or spin the stars up to corotation before merger~\cite{Kochanek1992,Bildsten1992} but those are not achievable for viscosities arising from nucleonic or leptonic weak processes. In particular, Lai (1994)~\cite{Lai_1994} found that shear and bulk viscous damping of mode oscillations of NS matter could only heat the stellar core to $\sim 10^8$K. Recent work by Arras et al. (2019)~\cite{Arras2019} considers heating due to direct Urca reactions during inspiral and concluded that they are heated upto $k_BT \sim 0.01$MeV which is one order of magnitude higher than  the estimates by Lai (1994)~\cite{Lai_1994}. But still these estimates show that NS remains cold during the inspiral and tidal heating does not produce any detectable effect in the inspiral GW signal.  \\

Hyperon bulk viscosity has been extensively studied over the past two decades mostly in the context of damping of $r-$mode oscillations from young, fast rotating pulsars~\cite{Jones2001,Lindblom2002,Debi_2006,Dalen2004,Jha2010,Haskell2010,Ofengeim,Alford2021,Jha2022}. It is found to strongly suppress the r-mode instability window below $10^{10}$K due to the high bulk viscosity produced by the non-leptonic weak-interaction channels involving hyperons~\cite{Lindblom2002,Haskell2010}. In this letter, we show that the high viscosity ($\approx 10^{8} - 10^{10}$ times the shear viscosity from $e-e$ scattering in the temperature region of $10^{6}-10^8$K) due to the presence of hyperons inside the high density core of the neutron star can sufficiently heat up the star during the inspiral to a much higher temperature estimated from earlier studies and these effects are potentially detectable as a deviation of the orbital decay rate from the general relativistic point-mass result using current and future generational GW detectors. 


\section{\label{sec:HyperonBV} 
Hyperonic Bulk viscosity}

%\dc{Given a strict page limit, this section can be shortened as most of this is available in cited literature.}\sg{Some sentence removed and some joined together, we are well within limits.}\\
 The EoS of nuclear matter inside NSs including hyperonic degrees of freedom can be described using different approaches~\cite{Chatterjee_2016,Oertel2017}. In this work we consider some standard parametrisations within the Relativistic Mean Field (RMF) model~\cite{GM91,Chen2014} where the strong interactions among baryons are mediated by mesons. Such EoS models have been widely applied to reproduce both nuclear matter properties as well as NS astrophysical data. Depending on the hyperon-nucleon potentials and the nuclear saturation parameters, hyperons appear at different onset densities inside the NS core. A detailed description of the chosen EoS parametrisation along with the hyperon potentials and onset densities of $\Lambda$ hyperons are given in Table~\ref{tab:table1}. \\

Bulk viscosity (BV) arises due to the pressure and density variations  associated  with the mode  oscillations inside the NS that  drive  the system out of chemical equilibrium~\cite{Jones2001}. 
%Strong  interaction  processes  do  not  play  a  role in producing bulk viscosity inside NSs,  because  the  strong  interaction  equilibrium  is  reached  so  fast that these processes can be considered to be in thermal equilibrium  compared  to  the  typical  pulsation  time  scales which is of the order of milliseconds~\cite{Dalen2004}. 
In addition to the leptonic hyperon processes, weak nonleptonic hyperon processes  contribute  more  efficiently to the bulk viscosity~\cite{Jones2001,Lindblom2002}. The standard non-leptonic reactions that contributes to the hyperon bulk viscosity are 
\begin{equation}\label{eq:Hyp1}
    \centering
    n + n \longleftrightarrow p + \Sigma^{-},
\end{equation}
\begin{equation}\label{eq:Hyp2}
    \centering
    n + p \longleftrightarrow p + \Lambda ,
\end{equation}
\begin{equation}\label{eq:Hyp3}
    \centering
    n + n \longleftrightarrow n + \Lambda .
\end{equation}
We do not consider the channel in Eq.~\eqref{eq:Hyp3} as this process has no simple $W$-exchange  contribution~\cite{Lindblom2002}. Among the nucleon-hyperon potentials, the best constrained potential is that of the hyperon $\Lambda$, having a value of $U_{\Lambda}=-30$ MeV ~\cite{Millener88,Schaffner92}. Although the potential depths for hyperons $\Sigma$ and $\Xi$ are not known precisely, it has been concluded that the $\Sigma$-nucleon potential is repulsive ~\cite{FriedmanGal07,SchaffnerGal00,Mares951} that restricts the appearance of $\Sigma$ hyperons inside NS matter. So, we only consider the reaction ~\eqref{eq:Hyp2} that contributes to the hyperon bulk viscosity. \\

For any weak interaction process, the real part of the bulk viscosity coefficient($\zeta$) is calculated in terms of relaxation times of microscopic processes
~\cite{Debi_2006}
\begin{equation}\label{BV}
    \zeta = n_B\frac{\partial P}{\partial n_n} \frac{dn_n}{d n_{B}}\frac{\tau}{1+(\omega\tau)^2},
\end{equation}
where  $P$ stands for the  pressure, $n_B$ is the total baryon number density , $n_n$ the neutron density, $\omega$  the oscillation frequency of the ($l,m$) mode and $\tau$ is the characteristic timescale of the reaction. The rates of these reactions can be calculated from the tree-level Feynman diagrams involving the exchange of a $W$-boson~\cite{Jones2001,Lindblom2002}.  For the weak non-leptonic  hyperon  processes  of Eqn.~\eqref{eq:Hyp2}, the  temperature ($T$) dependence of the relaxation time is given as $\tau \propto T^{-2}$~\cite{Lindblom2002}. From the typical expression of bulk viscosity, we see that it shows a resonant behaviour when the mode oscillation frequency ($\omega$) matches with the timescale ($\tau$). In Fig.~\ref{fig:fig_1}, we plot the bulk viscosity due to this weak-interaction channel for the chosen RMF parametrisations given in Table~\ref{tab:table1}, as a function of temperature for different mode frequency characteristic values; $1$ \& $2$ kHz for typical $f$-mode~\cite{Pradhan_2021} and $\sim 100$ Hz for $g$-mode oscillations~\cite{tran2022gmode}. Depending on the mode oscillation frequency, we see that the resonances occurs in the temperature range of $10^8 - 10^9$K and the resonant maximum value of $\zeta$ can reach a values of $10^{31} - 10^{32}$ gm cm$^{-1}$s$^{-1}$ which is several order of magnitude higher than the canonical shear and bulk viscosity ~\cite{Sawyer1989} in this temperature region. 
\begin{figure}
    \centering
    \includegraphics[width=0.5\textwidth]{Figures/Fig_1.pdf}
    \caption{Relative strengths of various sources of viscosities inside NS matter as a function of temperature for $n = 2.5n_0$( $n_0$ = the nuclear saturation density). The bands corresponds to hyperon bulk viscosities coming from  different EoSs given in Table~\ref{tab:table1} with the black solid line corresponding to HZTCS EoS. The standard neutron matter shear viscosity coming from $e-e$ scattering and bulk viscosity from m-URCA reactions at f = $1$kHz are also plotted along with their temperature dependence.}
    \label{fig:fig_1}
\end{figure}

\begin{table*}
\begin{adjustwidth}{-3em}{}
\begin{center}
    \begin{tabular}{|c|c|c|c|c|c|c|c|c|c|}
    \hline
    EOS & Max. mass & $\Lambda$ onset & Mass($M_{\odot}$) & Central  & Radius & Hyperon core& f-mode  & Max. temp(K)  & $\Delta \Phi = 2\pi \Delta\mathcal{N}$  \\ 
    & ($M_{\odot}$) & density & & density & (km) &  radius(km) & Freq.(Hz) &  at $D/R = 3$ & (rad) \\
    \hline
    \multicolumn{10}{c}{} \\
    \hline
    NL3 & 2.10 & 1.90$n_0$ & 1.6 & 2.07$n_0$ & 14.7 & 3.42 & 1847 & 9.7$\times 10^8$ & 0.001 \\
    ~\cite{Chen2014}& & & 1.8 & 2.48$n_0$ & 14.6 & 6.19 & 1909 & 3.3$\times 10^9$ & 0.08 \\
    & & & 2.0 & 3.35$n_0$ & 14.2 & 8.10 & 2009 & 6.2$\times 10^9$ & 0.3 \\
    \hline
    TM1 & 2.06 & 2.02$n_0$ & 1.6 & 2.24$n_0$ & 14.55 & 3.16 & 1873 & 8.7$\times 10^8$ & 0.0008 \\
    ~\cite{Steiner_2013}& & & 1.8 & 2.77$n_0$ & 14.37 & 6.18 & 1947 & 3.4$\times 10^9$ & 0.09 \\
    & & & 2.0 & 4.06$n_0$ & 13.6 & 8.15 & 2092 & 6.7$\times 10^9$ & 0.34 \\
    \hline
    TMA & 2.12 & 2.09$n_0$ & 1.8 & 2.54$n_0$ & 14.2 & 5.13 & 1948 & 2.3$\times 10^9$ & 0.02 \\
    ~\cite{Steiner_2013}& & & 2.0 & 3.35$n_0$ & 13.89 & 7.36 & 1909 & 5.1$\times 10^9$ & 0.16\\
    \hline
    HZTCS & 2.00 & 2.28$n_0$ & 1.6 & 2.67$n_0$ & 13.2 & 4.89 & 2108 & 2.3$\times 10^9$ & 0.02 \\
    ~\cite{HZTCS}& & & 1.8 & 3.32$n_0$ & 13.1 & 6.82 & 2171 & 4.7$\times 10^9$ & 0.16 \\
    & & & 2.0 & 5.32$n_0$ & 12.25 & 8.17 & 2305 & 7.9$\times 10^9$ & 0.44 \\
    \hline
    FSU2 & 2.03 & 1.92$n_0$ & 1.6 & 2.22$n_0$ & 14.4 & 4.98 & 1898 & 2.1$\times 10^9$ & 0.03 \\
    ~\cite{Chen2014}& & & 1.8 & 2.72$n_0$ & 14.2 & 7.07 & 1968 & 4.5$\times 10^9$ & 0.19 \\
    & & & 2.0 & 3.82$n_0$ & 13.6 & 8.54 & 2099 & 7.4$\times 10^9$ & 0.47 \\
    \hline 
    Stiffest EOS & 2.01 & 2.31$n_0$ & 1.6 & 2.71$n_0$ & 13.5 & 3.88 & 2047 & 1.4$\times 10^9$ & 0.004 \\
    from Ghosh et al.& & & 1.8 & 3.39$n_0$ & 13.4 & 6.34 & 2119 & 4.1$\times 10^9$ & 0.11 \\
    (2022)~\cite{Ghosh_2022}& & & 2.0 & 5.5$n_0$ & 12.5 & 8.05 & 2256 & 7.2$\times 10^9$ & 0.37 \\
    \hline
    \end{tabular} 
\end{center}
\end{adjustwidth}
    \caption{Detailed list of NS properties corresponding to different EOS models considered in this work. The potential depth of $\Lambda$ hyperons in normal nuclear matter is taken as $U_{\Lambda}^N = -30 $ MeV~\cite{Weissenborn_2012}. Densities are given in terms of $n_0$, the nuclear saturation density($\sim 0.15 $ fm$^{-3}$~\cite{Oertel2017}). The phase difference $\Delta\Phi$ at the last column denotes the total phase difference accumulated at the end of inspiral around $f \sim 500$Hz.}
    \label{tab:table1}
\end{table*}


\section{Tidal Heating in binary NS}
During the binary NS inspiral, the tidal energy transferred to the various oscillation modes will be damped out due to the viscous dissipation inside the NS matter. For normal nuclear matter inside NS in a binary, the main source of viscosity is the shear viscosity (SV) originating from electro-electron scattering in the relevant temperature region of $10^6 - 10^8$ K during the inspiral~\cite{Arras2019}. But in Fig.~\ref{fig:fig_1}, we showed that in this temperature range, hyperon bulk viscosity dominates the $e-e$ scattering shear viscosity by several orders of magnitude and also peaks at a temperature between $10^8 - 10^9$K depending on the mode oscillation frequency. To investigate the effect of this high viscous damping on the heating of the NS matter, let us first analyse the timescale for mode damping. The mode damping rate is given by
\begin{equation}\label{eq:damptime2}
    \gamma_{\alpha} = \Dot{E}_{visc,\alpha}/2E_{\alpha}
\end{equation}
where $\alpha$ denotes the particular eigenmode, $E_{\alpha}$ is the energy of the mode and $\Dot{E}_{visc,\alpha}$ is the energy dissipation rate. The rate of dissipated energy is given in terms of the viscous stress tensor $\sigma_{ij}$~\cite{Lai_1994}
\begin{equation}\label{eq:viscenergy}
    \Dot{E}_{visc} = \int d^3x \sigma_{ij}\bm{v}_{i,j}
\end{equation}
where $\bm{v}$ denotes the perturbation velocity vector and the viscous stress tensor $\sigma_{ij}$ can be written as~\cite{Landau1987Fluid}
\begin{equation}\label{eq:visctensor}
    \sigma_{ij} = \eta\left(\bm{v}_{i,j} + \bm{v}_{i,j} - \frac{2}{3}\delta_{ij}\nabla .\bm{v}\right) + \zeta\delta_{ij}\nabla . \bm{v}
\end{equation}
where $\eta$ and $\zeta$ are the shear and bulk viscosity coefficients respectively. \\

In the adiabatic approximation, the effect of the tidal potential due to the companion star is measured in terms of the Lagrangian fluid displacement vector $\bm{\xi}(r,t)$ from its equilibrium position. This displacement can be analysed in terms of the normal modes of the neutron star, 
\begin{equation}\label{eq:normal}
    \bm{\xi} (r,t) = \sum_{\alpha}\bm{\xi_{\alpha}}(r) e^{-i\omega_{\alpha}t}
\end{equation}
where $\alpha$ denotes the normal mode index and $\omega_{\alpha}$ is the angular frequency of the normal mode. For a particular mode, the eigenfunction can be written as a sum of the radial and tangential components 
\begin{equation}\label{eq:eigenvector}
    \bm{\xi_{\alpha}}(r) = \left[\zeta^r_{nl}(r)\bm{e_r} + r\zeta^{\perp}_{nl}(r)\bm{\nabla}\right]Y_{lm}(\theta,\phi)
\end{equation}
where $\bm{e_r}$ is the radial vector and $Y_{lm}(\theta,\phi)$ are the spherical harmonic functions. In this work, we will only consider the dominant fundamental ($f)$ mode contribution to the tidal heating. We determine the $f$-mode frequency via a relativistic Cowling approximation~\cite{Pradhan_2021} and the normalised mode eigenfunctions are used to calculate the tidal coupling~\cite{Andersson_2018}. \\

From the expression for the displacement vector $\xi$ given in Eq.~\eqref{eq:normal}, the velocity field is given by $\bm{v} = -i\omega \bm{\xi}$ (we drop the $\alpha$ subscript as we are working only with the $f$-mode). If, we consider only the bulk viscosity contribution to the energy dissipation given in Eqn.~\eqref{eq:viscenergy}, we can express the viscous dissipation rate as 
\begin{equation}\label{eq:gammabulk}
    \gamma_{bulk} = \frac{1}{2} \frac{(l+|m|)!}{(l-|m|)!}\int_0^R r^2dr\zeta \left(\frac{\partial \xi^r}{\partial r} + \frac{2}{r}\xi^r - l(l+1)\frac{\xi^{\perp}}{r}\right)^2 .
\end{equation}
We should compare this viscous dissipation timescale (which is given by inverse of the rate $\tau_{bulk} = 1/\gamma_{bulk}$) with the inspiral timescale to confirm whether the hyperon bulk viscous dissipation can effectively drain energy from the decaying orbit during the inspiral. We consider leading order gravitational radiation where the GW emission takes energy from the orbit at a rate of
\begin{equation}\label{eq:GW_energy}
    \Dot{E}_{gw} = \frac{-32\mathcal{M}\Omega}{5c^5}(G\mathcal{M}\Omega)^{7/3}~.
\end{equation}
Here the chirp mass is given by 
\begin{equation}\label{eq:chirp}
    \mathcal{M} = M\left(\frac{q^3}{1+q}\right)^{1/5}
\end{equation}
with the primary star mass is $M$ and the companion mass is  $qM$. $\Omega $ denotes the orbital frequency which is given by
\begin{equation}\label{eq:orbit}
    \Omega^2 = \frac{GM(1+q)}{D^3}
\end{equation}
where $D$ is the separation between the masses. In this Newtowian evolution dynamics, the orbital timescale is given as~\cite{Andersson_2018}
\begin{equation}\label{eq:time}
     \frac{1}{t_D} \equiv \frac{\Dot{\Omega}}{\Omega} = \frac{96}{5c^5}(G\mathcal{M}\Omega)^{5/3}\Omega~.
\end{equation}
\begin{figure}
    \centering
    \includegraphics[width=0.5\textwidth]{Figures/Fig_2.pdf}
    \caption{Estimated timescale for the different processes as a function of GW frequency compared against the inspiral timescale ($t_D$) for a NS binary of equal mass $1.8M_{\odot}$. Shear viscous(SV) and hyperon BV dissipation corresponds to the dominant $f$-mode dissipation by shear viscosity from $e-e$ scattering~\cite{Lai_1994} and BV from hyperons respectively. The tidal heating corresponds to the heating timescale from the hyperon BV dissipation. The different bands indicates uncertainities due to the choice of different EoSs given in Table~\ref{tab:table1}.}
    \label{fig:fig_2}
\end{figure}
In Fig.~\ref{fig:fig_2}, we plot this inspiral timescale for the relevant LIGO frequency band of $20-500$ Hz with the viscous dissipation and tidal heating timescale for the case of equal binary mass of 1.8M$_{\odot}$ for all the EOSs given in Table.~\ref{tab:table1}. We see that both the hyperon bulk viscous dissipation and tidal heating timescales are smaller than the inspiral timescale, confirming that unlike the case of shear viscous dissipation, the hyperon bulk viscous dissipation and heating happens faster than the orbital evolution and it can efficiently damp out the tidal energy  to heat up the star during the inspiral. \\

Now that we have established that bulk viscous dissipation will effectively damp out the mode oscillation during the inspiral, let us estimate the temperature the NS will reach till they come into contact. Since we are only considering the dominant mode $l = m = 2$ contribution, the energy dissipation rate can be estimated from the mode amplitude~\cite{Lai_1994}
\begin{equation}\label{eq:viscener}
    \Dot{E}_{visc} = \frac{12\pi}{5}\frac{GM^2}{R}q^2(1+q)\omega_0^{-4}Q_0^2\left(\frac{R}{D}\right)^92\gamma_{bulk}
\end{equation}
where $R$ is the radius of the star, $Q_0$  tidal coupling strength of the $f$-mode and $\omega_0$ the normalised frequency of the $f$-mode. The heat content of the star due to the degenerate fermionic gas in the core can be given as~\cite{Lai_1994}
\begin{equation}\label{eq:heat}
    U \approx 4.5\times 10^{45}T_8^2 erg = 4.5\times 10^{22}T^2 J .
\end{equation}
During the inspiral the thermal evolution of the star can be written as
\begin{equation}\label{eq:evol_T}
    \frac{dU}{dt} = \Dot{E}_{visc} + \Dot{E}_{cool},
\end{equation}
where $\Dot{E}_{cool}$ denotes the rate of cooling due to neutrino emission and surface photon luminosity. But the timescales for both these cooling processes are very high compared to the binary inspiral timescale as shown in Fig.~\ref{fig:fig_2}, and therefore can be neglected~\cite{Rees1992}. After integrating the thermal evolution equation~\eqref{eq:evol_T} from $D \to \infty$ when the stars were far apart and at a very low temperature ($10^5 - 10^6$K), we can get an estimate of the temperature reached as a function of their separation $D$
\begin{equation}\label{eq:final_T}
\begin{split}
    \left[\frac{T^4}{4} + Bln(T)\right] &=  \frac{\pi}{21870}\omega_0^{-4}Q_0^2q\frac{A}{10^{22}} \\ & \times \left(\frac{c^2R}{GM}\right)\left(\frac{c^3R^2}{G}\right)\left(\frac{3R}{D}\right)^5,
\end{split}    
\end{equation}
where $A$ and $B$ are parameters fitted to the functional dependence of $\gamma_{bulk}$ on the temperature ($T$), $\gamma_{bulk} = \frac{AT^2}{B+T^4}$ coming from the temperature dependence of timescale for hyperon BV. In Table~\ref{tab:table1}, we provide the estimates of the temperature reached at a separation of $D = 3R$ when the stars are about to merge. We see that, the temperatures are $\sim 10^9 - 10^{10}$K, which is twice order of magnitude higher than the earlier estimates~\cite{Lai_1994,Arras2019}. 

\section{Phase error estimation and detectability}
The energy loss due to tidal heating during the binary inspiral will lead to a change in the number of wave cycles ($\Delta\mathcal{N}$) or equivalently, a phase shift $\Delta\Phi = 2\pi\Delta\mathcal{N}$ in the observed frequency range of the GW detectors. This is very crucial to accurately guess the phase of the signal, otherwise the GW template can destroy a possible detection using matched filter technique~\cite{Matchfilter_2008}.  This additional torque to the viscous dissipation of energy will lead to a total change in the number of cycles given by
\begin{equation}\label{eq:N}
    \Delta\mathcal{N} = -\int_{f_a}^{f_b} t_D\left(\frac{\Dot{E}_{visc}}{\Dot{E}_{gw}}\right)df
\end{equation}
where $f_a$ is the frequency when the signal first enters the detector band, and $f_b$ , when it dives into the noise again. Since, $f$-modes are not resonantly excited, we need to do the integration over the whole frequency range unlike the cases for resonantly excited $g$-modes or $r$-modes where additional energy loss is associated with the particular mode frequency~\cite{Andersson_2018,Ma2021}. Using the expressions given in Eqn.~\eqref{eq:GW_energy} and ~\eqref{eq:viscener}, the net change in the number of cycles is given by
\begin{equation}\label{eq:error}
   \Delta\mathcal{N}  = \frac{15\pi}{384}\left(\frac{c^2R}{GM}\right)^6 \omega_0^{-4}Q_0^2\frac{1}{q(1+q)}\left(\frac{R}{c}\right)^2\int_{f_a}^{f_b} \gamma_{bulk} df.    
\end{equation}
%Since, $\gamma_{bulk}$ has a temperature dependence which also change as a function of frequency (or equivalently the separation $D$) according to Eqn.~\eqref{eq:final_T}.
In Table~\ref{tab:table1}, we demonstrate the net phase difference accumulated at GW frequency of $500$ Hz for all the different EoSs and different values of equal mass binaries. In Fig.~\ref{fig:fig_3}, we display how this phase difference grows as a function of frequency, taking into consideration uncertainties due to the different choice of EOSs. For equal $1.6M_{\odot}$ binaries, we see that the net phase difference is of the order of $10^{-3} - 10^{-2}$ rad and for higher masses of $1.8M_{\odot}$ or $2M_{\odot}$,  we get a net phase difference in the order of $0.1-0.5$ rad. \\

To be able to measure this phase difference using the current or future generation GW detectors, the phase uncertainty of detected GWs must be smaller than the phase shift. Earlier estimates based on a detection with signal-to-noise ratio (SNR) of 10 and an approximate single detector sensitivity  showed that phase uncertainty was to be around $\Delta\Phi \approx 1 - 3$ rad ($\Delta\mathcal{N}$ = 0.5 equivalent to $\Delta\Phi = \pi$)~\cite{Cutler_1994}. However recent improvements of analysis~\cite{Bruce_2021} have shown the phase error to be around $\Delta\phi \sim \pm 0.1$ rad for $f_{GW} \leq 300$ Hz inclusive of calibration uncertainties  for the GW170817 signal analysed using GW waveform
model \textit{IMRPhenomPv2\_NRTidal}~\cite{NRTidal}. More recently, Read (2023)~\cite{read2023waveform} compares a number of GW waveform models and shows that the uncertainty due to waveform differences is $\sim \pm 0.02$ rad for A+~\cite{ALIGO} and $\pm 10^{-3}$ rad for Cosmic Explorer (CE)~\cite{CE}. So, from these estimates, we see that a BNS event with SNR like GW170817 would produce enough tidal heating to be detectable using the current LVK detectors if it has a heavier component mass $\geq 1.8M_{\odot}$. In the 3G detectors, we can measure evidence of tidal heating due to hyperons even for much lower mass NS components. 
\begin{figure}
    \centering
    \includegraphics[width=0.5\textwidth]{Figures/Fig_3.pdf}
    \caption{Estimated phase shift accumulated in the GW signal for equal mass binaries as a function of GW frequency. The different bands indicate uncertainties due to the choice of different EoSs given in Table~\ref{tab:table1} corresponding to different masses of equal mass binaries considered and the solid lines define their boundaries.}
    \label{fig:fig_3}
\end{figure}

\section{Discussion} 
This work considers for the first time the effect of tidal heating in NSs during binary inspirals due to bulk viscosity originating from non-leptonic weak interaction processes involving hyperons inside the NS core. We consider several state-of-the-art EoSs including the hyperonic degrees of freedom consistent with multi-messenger observations and calculate the bulk viscous dissipation of the dominant $f-$mode oscillations excited due to the tidal interaction during the inspiral. This dissipated energy can be effectively converted to thermal energy during the inspiral timescale and can heat up the stars upto $0.1 - 1$ MeV during the last orbits before coalescing.
\\

For systems with mass ratio $q>1$ such as neutron star-black hole (NSBH) systems, it is evident from Eqn.~\eqref{eq:final_T} and ~\eqref{eq:error} that the temperature estimate will be higher by a factor of $q^{1/4}$ and the net phase shift decreases by a factor of $q(1+q)$ respectively. Although these estimates are higher by orders of magnitudes than earlier estimates by Lai (1994)~\cite{Lai_1994} and Arras et al. (2019)~\cite{Arras2019}, they are not sufficiently high to demand inclusion of thermal effects in the EoS during inspiral. Recent studies~\cite{Thermal_2022} have shown that thermal correction to the tidal deformability and radius is negligible for temperatures below $1$ MeV.  \\

 For neutron stars with lower mass $\sim1.4M_{\odot}$, the central density may be too low support sufficient hyperon fraction to produce significant tidal heating. Recent analysis  of populations of galactic binary NSs by Farrow et al.(2019)~\cite{Farrow_2019} suggests two distinct mass distributions for recycled and slow NSs, and bimodality of the recycled NS mass distribution. Although this model predicts that we need $\mathcal{O}(10)$ and $\mathcal{O}(100)$ BNS events to have a component mass $\geq 1.6M_{\odot}$ and $\geq 1.8M_{\odot}$ respectively, recent observations of BNS event GW190425~\cite{GW190425} or NSBH event GW200105~\cite{NSBh} having component NS masses $m_1^{90\%} = [1.6,1.87]M_{\odot}$ and $m_2 = 1.9^{+0.3}_{-0.2}$ respectively, suggest that there might be other formation channels for higher mass NS in binary systems. The induced phase shift to the GW signal due to the tidal heating is $\sim 0.1$ rad for component masses $\geq 1.8M_{\odot}$ which is detectable with current and future generation GW detectors, although a detailed post-Newtonian (PN) analysis of the effect of viscous dissipation on the GW phasing is required to accurately determine bulk viscosity from a future detection. A recent work by Most et al.~(2022)~\cite{Most_2021} performed an order-of-magnitude PN estimate about the direct effect of bulk viscosity on the orbital motion ignoring back-reaction of viscosity on the mode oscillations that leads to tidal heating. Considering their estimates, at $f = 100$Hz, when the bulk viscosity is $\sim 10^{32}$ gm cm$^{-1}$ s$^{-1}$(from Fig.~\ref{fig:fig_1}), we would require a SNR of $100$ to get a $10\%$ measurement on the bulk viscosity which is achievable using the third generation GW detectors. A future detection of this high bulk viscosity effect originating from non-leptonic weak interaction processes involving hyperons either in terms of deviation in the orbital motion or tidal heating would unequivocally confirm the presence of hyperons in the interior of neutron stars. Even a non-detection can place an upper limit on the bulk viscosity and the hyperon fraction inside the neutron star which is also very crucial for  modelling of dense nuclear matter. \\

There are various directions in which our research can be further developed. First, a detailed investigation of the PN order at which the tidal heating becomes relevant will have to be performed to confirm the conclusions related to the qualitative Newtonian inspiral discussed in this work. Second, we only considered the contribution of the dominant $f$-mode dissipation but Lai (1994)~\cite{Lai_1994} showed that tidal heating due to the resonant $g-$mode oscillations can also be as significant in spite of their tidal coupling coefficient being too small ($\sim 10^{-3}-10^{-4}$~\cite{Andersson_2018}). Due to their frequency being $\sim 100$ Hz, the bulk viscosity value is also one order of magnitude order higher than those of $f-$modes, so the tidal heating due to resonant $g$-modes should also be studied in detail. Third, the dominant channel in unpaired strange quark matter, $d + u \longleftrightarrow s + u$ process, can also produce bulk viscosity as high as $10^{29}-10^{30}$ gm cm$^{-1}$ s$^{-1}$ in the temperature region of $10^8-10^9$K~\cite{Alford_2010,Alford_2014}. Although this estimate is one order of magnitude less than those of hyperons, it is still worthwhile to see the amount of tidal heating produced in a strange star during binary inspiral. Finally, we need to consider the effect of superfluidity on the hyperon bulk viscosity since they are known to reduce the rates of the weak interactions and also effective below the critical temperature of $10^9$K~\cite{Haskell2010}. These topics are currently work in progress and will be reported in forthcoming publications. \\


\begin{acknowledgments}
 The authors would like to acknowledge the usage of the IUCAA HPC computing facility, PEGASUS for the numerical calculations. 
\end{acknowledgments}



\nocite{*}

\bibliography{apssamp}% Produces the bibliography via BibTeX.

\end{document}
%
% ****** End of file apssamp.tex ******
